\section{Conclusion}

This paper establishes the Neural Tangent Kernel (NTK) theory for low-rank random feature architectures (MMNN/RF-LR), demonstrating that these models provide a \textbf{clean and advantageous route} to the kernel regime with improved theoretical guarantees and practical efficiency.

Our primary contribution shows that there is \textbf{no RKHS disadvantage} for low-rank random features relative to standard fully-connected networks when the rank \( r \) is moderately large: the induced RKHS coincides in the relevant limits, so expressivity is preserved. More importantly, we prove that the \textbf{kernel regime is accessible at a linear parameter budget \( O(rn) \)} rather than the quadratic \( O(n^2) \) of dense layers. This \textbf{linear-width scaling} is not merely computational—it \emph{lowers the threshold} at which NTK predictions become accurate, enabling theoretical guarantees at markedly smaller model sizes than previously possible.

Beyond expressivity equivalence, we establish improved concentration properties. In a doubly-random Gram matrix regime, NTK fluctuations decay as \textbf{\( O(1/r^2) \)}, which is \emph{significantly faster} than the classical \textbf{\( O(1/N) \)} decay from standard NTK theory. For fixed total parameters \( N_{\text{total}} = rn \), this yields a \textbf{quadratic reliability gain} for kernel approximations when \( r \) is moderate and \( n \) is large. Practically, this means the \textbf{deterministic NTK limit is reached sooner, with smaller models}, which is a core theoretical advantage of the low-rank random feature design in the lazy regime.

The spectral analysis at initialization reveals an informative \textbf{bulk–outlier structure} with a \textbf{single dominant outlier} atop a bulk consistent with \textbf{Marchenko–Pastur} behavior, aligned with Terj\'ek-style analyses. We show that the outlier controls the fastest decaying direction under kernel gradient flow, while the ill-conditioned bulk explains slow directions and the \textbf{bad conditioning observed at initialization}. The low-rank structure renders this outlier analysis more tractable and facilitates cleaner comparisons to fully-connected networks at \textbf{matched parameter budgets}.

We develop a framework for finite-width corrections that organizes deviations from the infinite-width limit into a \textbf{\( 1/r \) expansion across depth \( L \)}, with explicit \textbf{\( O(1/r) \)} and \textbf{\( O(1/r^2) \)} effects. Two-layer configurations are \textbf{especially tractable}, allowing precise cross-checks between theory, empirical NTKs, and training trajectories. We further identify \textbf{corrections scaling in \( Nr \)} due to edge-of-chaos considerations for MMNNs, providing a concrete route to quantify departures from the kernel limit in realistic models.

From a practical standpoint, our analysis suggests \textbf{\( r \in [5,50] \)} as a regime that balances concentration, expressivity, and parameter efficiency across tasks. Sample complexity heuristics align with the spectrum-of-kernel-random-matrices viewpoint, pointing to \textbf{linear scaling in ambient dimension} in appropriate regimes. We clarify how \textbf{data normalization matters}: on the sphere, \textbf{\( \operatorname{Tr}=1 \)} removes certain \( 1/r \) normalization effects and clarifies which NTK variant is appropriate; in hypercube settings, \textbf{\( \operatorname{Tr}(\Sigma_w)=d \)} with unit bias preserves approximation conditions. These choices \textbf{directly affect NTK variance} and should be matched to the data domain.

We analyze the architectural policy of training only output weights \( A^{(\ell)} \) while freezing input weights \( w^{(\ell)} \), showing that this \textbf{partial training reduces the NTK magnitude by approximately a factor of two} relative to fully trainable counterparts. In the kernel regression regime, this \textbf{effectively doubles the learning-rate scale} for the same optimizer hyperparameters (equivalently, halves the kernel eigenvalues), modifying convergence speed while preserving the RKHS solution.

In summary, viewed strictly through the NTK lens, low-rank random features provide a \textbf{clean and advantageous route} to the kernel regime: \textbf{the same RKHS} as dense networks, \textbf{entry at linear budgets}, and \textbf{faster concentration \( O(1/r^2) \)} than classical \( O(1/N) \) theory predicts. The \textbf{bulk–outlier spectral picture} illuminates early training and conditioning. These results establish a solid theoretical foundation for understanding MMNNs in the lazy training regime and provide concrete guidance for practical deployment.

\section{Discussion}

While this paper establishes core NTK results for low-rank random feature architectures, several directions remain open and define an immediate research roadmap. We hope to develop a systematic \textbf{finite-width corrections} program for this architecture to elucidate the \textbf{finite-width loss landscape} under the NTK perspective specifically for low-rank random-feature models, building on the \( 1/r \) expansion framework introduced here.

Several \emph{NTK-specific} theoretical questions are marked for future work. A \textbf{non-Gaussian process propagation analysis} distinct from NNGP fits would provide complementary insights beyond the kernel regime. An \textbf{explicit bulk formula} with normalization scaled by \textbf{\( r^2 \)} remains to be fully characterized. A \textbf{concentration bound for \( 2 \times \text{NTK}_{\text{MMNN}} - \text{NTK}_{\text{MLP}} \)} at large \( r \) would formalize the relationship between low-rank and dense architectures. Several theorem-level items require completion: \textbf{concentration bounds} and \textbf{outlier analyses} marked as ``to be filled,'' remaining \textbf{finite-width NTK} parts and their relation to \textbf{\( Nr \)} scaling, and a \textbf{formal proof} of the aforementioned \textbf{factor-of-two} NTK magnitude reduction under partial training.

Additional theoretical extensions include investigating \textbf{Stiefel-manifold} assumptions and their implications in the extensive-rank regime, developing \textbf{DSRN} (Deep Structured Random Networks) extensions of NTK theory, and exploring how the curse of dimensionality impacts normalization choices (unit norm weights vs. biases) depending on data geometry (hypercube vs. spherical).

The experimental roadmap to validate these theoretical predictions is equally concrete. Future work should include \textbf{confirming lazy-kernel predictions} and identifying the boundary of lazy training in practice; \textbf{empirically validating \( r \in [5,50] \)} as a strong-concentration window across diverse tasks; conducting \textbf{two-layer width sweeps} to test \( O(1/n)+O(1/r) \) finite-width corrections and identify minimal \( (n,r) \) configurations for accurate kernel predictions; and a comprehensive \textbf{comparison suite with MLPs} to probe NTK spectra, staircase loss dynamics on multi-index targets, Hessian/condition-number evolution through training, and task accuracy at matched parameter counts. These studies are designed to verify the \textbf{\( O(1/r^2) \) concentration}, the \textbf{single-outlier structure}, and the \textbf{practical accessibility of the kernel regime}.

Our scope in this work is limited to the \textbf{lazy/NTK regime} for the low-rank RF architecture. The broader mean-field picture (stepwise dynamics, central-flow arguments, dynamical stability) belongs to a distinct regime and constitutes a separate research program. Completing the pending finite-width formulas, spectrum characterizations, and matched-budget experiments will turn this into a comprehensive NTK account of MMNNs and lay a solid foundation for analyses beyond the lazy regime, including connections to mean-field dynamics and global convergence theory.
